This appendix expands the consolidated downstream results (\Cref{tab:bench}) into the per-component, per-loss breakdown from our evaluation round, and adds two measurements the main table omits: general visual question answering (GQA) and a second counting benchmark (TallyQA, FSC-147). All numbers are zero-shot.

\paragraph{Capability preservation.}
\Cref{tab:gqa_pope} reports GQA and POPE. The picture sharpens the dissociation of \Cref{sec:analysis-pheno}: aligning the \emph{language model} (alone or with the projector) leaves general VQA and hallucination essentially untouched (POPE F1 even ticks up), whereas aligning the \emph{projector alone} degrades both markedly (GQA $0.606\!\to\!0.580$, POPE F1 $0.839\!\to\!0.721$). This is consistent with the drift analysis (\Cref{sec:analysis-where}): the projector-only run moves the projector hardest, feeding the frozen language model visual features that have drifted off its training distribution.

\begin{table}[h]
\centering
\caption{General VQA and hallucination (zero-shot). Higher is better; best per column in \textbf{bold}. Aligning the language model preserves capability; aligning the projector alone degrades it.}
\label{tab:gqa_pope}
\vskip 0.1in
\small
\begin{tabular}{lccc}
\toprule
Variant & GQA acc. & POPE acc. & POPE F1 \\
\midrule
LLaVA-1.5-7B (base)        & \textbf{0.606} & 0.854 & 0.839 \\
Caption SFT ($\lambda{=}0$) & 0.602 & 0.857 & 0.844 \\
KL align, projector only   & 0.580 & 0.778 & 0.721 \\
KL align, LM only          & 0.596 & \textbf{0.859} & \textbf{0.846} \\
KL align, LM + projector   & 0.599 & 0.852 & 0.836 \\
\bottomrule
\end{tabular}
\end{table}

\paragraph{Counting.}
\Cref{tab:counting} reports TallyQA accuracy as a function of which component is trained and whether the alignment loss is on. The pattern is striking: training the \emph{projector} improves counting (even without the loss, $40.96\!\to\!42.67$; with it, $45.06$), while training the \emph{language model} degrades it ($\to\!35.82$ without the loss, $34.38$ with), and the joint run follows the language model down. The alignment loss does not create these effects so much as \emph{amplify} each component's native direction.

Crucially, language-model alignment is exactly the configuration that produces the largest saliency gains (\Cref{tab:align}) yet the worst counting, so across components, saliency-map quality and counting accuracy are not merely uncorrelated but \emph{anti}-correlated. This is the same message as the attention-overlap-vs-counting null of \Cref{app:counting}, now seen across training configurations. (Counting benchmarks themselves disagree on the projector, which is best here on TallyQA but worst on CountBench in \Cref{tab:bench}; the robust signal is that language-model alignment never improves counting despite producing the best maps.) FSC-147 is near chance for every variant ($\approx 5$--$6\%$ accuracy, with diverging mean error), so we omit it from the table; the models simply cannot do open-set counting at these magnitudes.

\begin{table}[h]
\centering
\caption{TallyQA counting accuracy (\%) by trained component and loss. Base (no fine-tuning) scores $40.96$. Training the projector helps counting; training the language model hurts it; the loss amplifies the component's direction.}
\label{tab:counting}
\vskip 0.1in
\small
\begin{tabular}{lcc}
\toprule
Trained component & SFT ($\lambda{=}0$) & + KL align ($\lambda{=}0.5$) \\
\midrule
Projector            & 42.67 & \textbf{45.06} \\
Language model       & 35.82 & 34.38 \\
LM + projector       & 37.25 & 33.57 \\
\bottomrule
\end{tabular}
\end{table}
